\documentclass[11pt]{article}

\usepackage[T1]{fontenc}
\usepackage{amsmath,amssymb}
\usepackage{booktabs}
\usepackage{hyperref}

\title{Classical Mechanics as an Explicit Set-Theoretic Predicate}
\author{Formalization notes}
\date{\today}

\begin{document}
\maketitle

\begin{abstract}
This note expresses classical particle mechanics in explicit predicate form,
emphasizing the biconditional structure
$\operatorname{CM}(\mathcal{M}) \iff P_1 \land P_2 \land \cdots \land P_6$.
The formulation clarifies the direct equivalence between the definition of
a classical mechanics model and the conjunction of its constitutive axioms.
A Lean 4 encoding demonstrates how the predicate selects exactly those
systems satisfying all six conditions, and how incompatibility arises when
conflicting axiomatizations are present.
\end{abstract}

\section{The Classical Mechanics Predicate}

We define the Classical Mechanics predicate as a biconditional over a system
$\mathcal{M} = (P, T, m, s, f)$ with signature components and auxiliary
structure as in the foundational formulation. The predicate is:

\[
\operatorname{CM}(\mathcal{M}) \quad\Longleftrightarrow\quad
  P_1 \land P_2 \land P_3 \land P_4 \land P_5 \land P_6,
\]

where the six conditions are:
\begin{enumerate}
  \item[\(P_1\)] The particle set $P$ is finite and nonempty.
  \item[\(P_2\)] The time domain $T$ is an interval in $\mathbb{R}$.
  \item[\(P_3\)] For each particle $p \in P$, position $s(p, \cdot)$ is
    twice differentiable on $T$.
  \item[\(P_4\)] For each particle $p \in P$, mass $m(p) > 0$.
  \item[\(P_5\)] For each pair $(p, t) \in P \times T$, the force series
    $\sum_{i \in I} f(p, t, i)$ converges absolutely.
  \item[\(P_6\)] For each pair $(p, t) \in P \times T$, Newton's second law holds:
    \[
      \sum_{i \in I} f(p, t, i) = m(p)\, \frac{\partial^2 s}{\partial t^2}(p, t).
    \]
\end{enumerate}

\subsection{Predicate as Logical Equivalence}

The predicate $\operatorname{CM}(\mathcal{M})$ is defined such that:
\[
  \mathcal{M} \text{ is a classical mechanics model} \quad\Longleftrightarrow\quad
  \mathcal{M} \text{ satisfies } P_1 \land P_2 \land \cdots \land P_6.
\]

This formulation makes explicit that to verify membership in the class of
classical mechanics models, one must check all six conditions, and the
conjunction is both necessary and sufficient.

\section{Lean Encoding}

The Lean module \texttt{ClassicalMechanicsPredicate.lean} encodes this
predicate-based approach. The core definition is:

\begin{verbatim}
def isClassicalMechanicsModel (S : Signature) (M : System S) : Prop :=
  (Nonempty S.Particle) ∧
  (HasFiniteCarrier S.Particle) ∧
  (S.realTimeInterval) ∧
  (∀ p, S.twiceDifferentiable (M.position p) (M.acceleration p)) ∧
  (∀ p, S.positive (M.mass p)) ∧
  (∀ p t, S.absolutelySummable (M.force p t)) ∧
  (∀ p t, S.forceSum (M.force p t) =
    S.scale (M.mass p) (M.acceleration p t))
\end{verbatim}

The theorem \texttt{cm\_iff\_conjunction} establishes the equivalence by
reflexivity, and \texttt{system\_satisfies\_predicate} shows that any
\texttt{System} structure satisfying the constructor constraints makes
\texttt{isClassicalMechanicsModel} true.

\section{Conflict Detection via Incompatible Axiomatizations}

When a second module \texttt{FakeSciencePredicate.lean} is introduced that
axiomatizes conflicting requirements (e.g., declaring a false axiom in the
same namespace), the Lean type checker detects the incompatibility. The module
\texttt{PredicateAudit.lean} attempts to combine both formulations, which
reveals:

\begin{enumerate}
  \item Redeclaration conflict: identifiers cannot be declared twice in the
    same namespace.
  \item False axiom dependency: any theorem derived from the conflicting module
    depends on a false axiom, violating consistency.
  \item Type-checking failure: Lean halts compilation and reports the specific
    conflict location.
\end{enumerate}

This demonstrates how the predicate-based approach, once formalized, provides
automatic verification of consistency: if two axiomatizations cannot coexist
in a single Lean environment, they are logically incompatible.

\section{Philosophical Implications}

In the Suppes tradition, a set-theoretic predicate identifies a class of
mathematical structures that are models of a theory. The explicit biconditional
form $\operatorname{CM}(\mathcal{M}) \iff \bigwedge P_i$ emphasizes that:

\begin{enumerate}
  \item The predicate is a pure \textit{definition}, not an empirical claim.
  \item A structure is a classical mechanics model if and only if it satisfies
    all six conditions; there is no partial membership.
  \item The Lean formalization serves as a certificate of structural coherence:
    if the module type-checks, the predicate is internally consistent.
  \item Incompatibility between rival theories is \textit{decidable} via the
    module system: conflicts are caught at compile time, not deferred to
    semantic analysis.
\end{enumerate}

\section{Comparison with Foundational Formulation}

The foundational formulation (\texttt{ClassicalMechanics.lean}) uses a
\texttt{System} structure that bundles the conditions as constructor fields,
ensuring that only systems meeting all requirements can be instantiated. The
predicate formulation (\texttt{ClassicalMechanicsPredicate.lean}) instead
provides an explicit test function that returns true precisely when all
conditions hold. Both approaches are equivalent in power, but the latter
makes the conjunctive structure and the biconditional semantics transparent.

\bibliographystyle{plain}
\bibliography{references}

\end{document}
